\documentclass[11pt]{article}
\usepackage[margin=1in]{geometry}
\usepackage{amsmath, amssymb, amsthm, mathtools}
\usepackage{hyperref}

\newtheorem{theorem}{Theorem}
\newtheorem{lemma}[theorem]{Lemma}
\theoremstyle{definition}
\newtheorem{definition}[theorem]{Definition}
\theoremstyle{remark}
\newtheorem*{remark}{Remark}

\title{Optimal Transport and the Wasserstein Distance}
\author{math-foundations \textbullet{} concept 41}
\date{}

\begin{document}
\maketitle

\subsection*{First, an example}

Take two point masses: all probability mass sitting at $0$ for $\mu = \delta_0$, and
all mass sitting at $1$ for $\nu = \delta_1$.  To turn $\mu$ into $\nu$ you must
physically move the unit of mass a distance $1$, so
\[
W_1(\mu, \nu) \;=\; 1.
\]
\emph{Read aloud:} ``W-one of mu and nu equals one.''
By contrast, $D_{\mathrm{KL}}(\mu \,\|\, \nu) = +\infty$ because the supports do not
overlap; KL cannot tell us \emph{how far} the mass moved, only that it moved.
For translated unit Gaussians, $W_1(\mathcal{N}(0,1), \mathcal{N}(m,1)) = |m|$ grows
linearly in the shift while $D_{\mathrm{KL}} = m^2/2$ grows quadratically.
Wasserstein is the distance for distributions that move \emph{in space}, not in
shape.

\section*{Setting}

Let $(\mathcal{X}, d)$ be a Polish metric space (e.g.\ $\mathbb{R}^n$ with the Euclidean
metric) and let $\mathcal{P}_p(\mathcal{X})$ denote the set of Borel probability
measures on $\mathcal{X}$ with finite $p$-th moment, for some $p \geq 1$.

\begin{definition}[Coupling / transport plan]
Given $\mu, \nu \in \mathcal{P}_p(\mathcal{X})$, a \emph{coupling} of $\mu$ and $\nu$
is a Borel probability measure $\pi$ on $\mathcal{X} \times \mathcal{X}$ whose marginals
are $\mu$ and $\nu$:
\[
\pi(A \times \mathcal{X}) = \mu(A), \qquad \pi(\mathcal{X} \times B) = \nu(B).
\]
\emph{Read aloud:} ``pi of A cross X equals mu of A, and pi of X cross B equals nu of B.''
We write $\pi \in \Pi(\mu, \nu)$.  The set $\Pi(\mu, \nu)$ is always non-empty
(contains $\mu \otimes \nu$) and weak-$*$ compact.
\end{definition}

\begin{definition}[Wasserstein-$p$ distance, Kantorovich formulation]
\[
W_p(\mu, \nu) \;=\; \left( \inf_{\pi \in \Pi(\mu, \nu)}
   \int_{\mathcal{X} \times \mathcal{X}} d(x,y)^p \, d\pi(x,y) \right)^{1/p}.
\]
\emph{Read aloud:} ``W-p of mu and nu equals the infimum over couplings pi of the integral of d-of-x-y to the p-th power against pi, all to the one-over-p.''
By weak-$*$ compactness of $\Pi(\mu, \nu)$ and lower semi-continuity of the cost,
the infimum is attained.
\end{definition}

\begin{definition}[Monge formulation]
A \emph{transport map} $T : \mathcal{X} \to \mathcal{X}$ pushes $\mu$ forward to $\nu$
if $T_\#\mu = \nu$, i.e.\ $\nu(B) = \mu(T^{-1}(B))$ for every Borel $B$.  The Monge
problem is
\[
\inf_{T_\#\mu = \nu} \int_{\mathcal{X}} d(x, T(x))^p \, d\mu(x).
\]
\emph{Read aloud:} ``infimum over transport maps T pushing mu to nu of the integral of d-of-x and T-of-x to the p, against mu.''
Every map $T$ induces a (deterministic) coupling $(\mathrm{id}, T)_\# \mu$, so the
Monge value is always $\geq$ the Kantorovich value.  Under regularity
(e.g.\ $\mu$ absolutely continuous on $\mathbb{R}^n$, $p > 1$), Brenier's theorem
gives equality and a unique optimal map.
\end{definition}

\section*{The 1D case: a closed-form theorem}

In one dimension everything collapses to ordering quantiles.  Recall that for a
probability measure $\mu$ on $\mathbb{R}$ with c.d.f.\ $F_\mu(x) = \mu((-\infty, x])$,
the \emph{quantile function} is
\[
F_\mu^{-1}(u) \;=\; \inf\{ x \in \mathbb{R} \;:\; F_\mu(x) \geq u \}, \qquad u \in (0,1).
\]
\emph{Read aloud:} ``F-mu inverse of u equals the infimum of those real x for which F-mu of x is at least u, for u between zero and one.''
If $U \sim \mathrm{Unif}(0,1)$ then $F_\mu^{-1}(U) \sim \mu$.

\begin{theorem}[1D Wasserstein closed form]
\label{thm:1d}
Let $\mu, \nu \in \mathcal{P}_p(\mathbb{R})$ with $p \geq 1$.  Then
\[
W_p(\mu, \nu)^p \;=\; \int_0^1 \left| F_\mu^{-1}(u) - F_\nu^{-1}(u) \right|^p \, du.
\]
\emph{Read aloud:} ``W-p of mu and nu, raised to the p, equals the integral from zero to one of the absolute difference of the two quantile functions, raised to the p, du.''
Moreover the infimum is attained by the \emph{monotone coupling}
$\pi^* = (F_\mu^{-1}, F_\nu^{-1})_\# \mathrm{Leb}_{(0,1)}$.
\end{theorem}

\begin{proof}[Proof sketch]
\textbf{Upper bound.}  The map $u \mapsto (F_\mu^{-1}(u), F_\nu^{-1}(u))$ from
$(0,1)$ to $\mathbb{R}^2$ pushes Lebesgue measure on $(0,1)$ to a probability $\pi^*$ on
$\mathbb{R}^2$.  Its first marginal is the law of $F_\mu^{-1}(U) = \mu$, and similarly
for $\nu$, so $\pi^* \in \Pi(\mu, \nu)$.  Hence
\[
W_p(\mu, \nu)^p \;\leq\; \int_{\mathbb{R}^2} |x-y|^p \, d\pi^*(x,y)
                  \;=\; \int_0^1 |F_\mu^{-1}(u) - F_\nu^{-1}(u)|^p \, du.
\]
\emph{Read aloud:} ``W-p to the p is at most the integral of x-minus-y to the p against pi-star, which equals the quantile-difference integral from zero to one.''

\textbf{Lower bound (the key step).}  The cost $c(x,y) = |x-y|^p$ for $p \geq 1$
satisfies the \emph{Monge condition} (a.k.a.\ submodularity):
for all $x < x'$ and $y < y'$,
\[
c(x, y) + c(x', y') \;\leq\; c(x, y') + c(x', y).
\]
\emph{Read aloud:} ``c-of-x-y plus c-of-x-prime-y-prime is less than or equal to c-of-x-y-prime plus c-of-x-prime-y.''
(For $p = 1$ this is immediate: both sides equal $|x - y'| + |x' - y|$ in the
``crossing'' configuration but $(x'-x) + (y'-y)$ vanishes when one rearranges; the
general case follows from convexity of $t \mapsto |t|^p$.)

Now let $\pi$ be any coupling and consider any pair of points
$(x, y), (x', y') \in \mathrm{supp}(\pi)$ with $x < x'$ and $y > y'$.  The Monge
condition says we can ``uncross'' this pair: replace mass $\varepsilon > 0$ on these
two pairs with mass $\varepsilon$ on $(x, y')$ and $(x', y)$ instead, without changing
either marginal, and the total cost can only decrease.  An exchange/measurable-
selection argument (a swapping/cyclical-monotonicity lemma; see Villani Ch.~5)
shows that any coupling can be made \emph{monotone} (i.e.\ supported on a non-
decreasing set) without increasing cost.  But the only monotone coupling with the
prescribed marginals is $\pi^*$.  Hence $\pi^*$ is optimal, proving the lower
bound.
\end{proof}

\begin{remark}
The theorem yields the trivial empirical-distribution algorithm: to estimate
$W_p$ between samples $x_1, \ldots, x_n$ and $y_1, \ldots, y_n$, sort each list
and compute $\bigl( \tfrac{1}{n} \sum_i |x_{(i)} - y_{(i)}|^p \bigr)^{1/p}$.
\end{remark}

\section*{Examples}

\textbf{Two point masses.}  $W_p(\delta_a, \delta_b) = |a - b|$ for every $p$.

\smallskip

\textbf{Translated Gaussians.}  $\mu = \mathcal{N}(0, 1), \nu = \mathcal{N}(m, 1)$:
$W_p(\mu, \nu) = |m|$.  Compare with $D_{\mathrm{KL}}(\mu \,\|\, \nu) = m^2/2$.

\section*{Pointers}

The Kantorovich--Rubinstein duality
$W_1(\mu, \nu) = \sup_{\|f\|_{\mathrm{Lip}} \leq 1} \int f \, d(\mu - \nu)$ is the
gateway to Wasserstein GANs.  Brenier's theorem identifies the optimal transport
map for $p = 2$ on $\mathbb{R}^n$ as the gradient of a convex function.  See Villani's
\emph{Optimal Transport: Old and New} and Peyr\'e--Cuturi's
\emph{Computational Optimal Transport}.

\end{document}
